\subsection{Abstand eines Punktes von einer Ebene}
\label{subsec:analytic_geometry_point_to_plane_distance}

Bei einer Ebene zeigt der Normalenvektor bereits in Richtung der kürzesten
Verbindung. Deshalb genügt eine skalare Projektion auf diese Richtung.

\begin{center}
  \CachedTikzFigure{content/resources/chapters/analytic_geometry/tikz/point_plane_distance.tex}
\end{center}

\begin{TheoremBox}{Punkt--Ebenen-Abstand}{analytic_geometry_point_plane_distance}
  Sei $P$ ein Punkt der Ebene $E$ mit Ortsvektor $\vec p$, sei $\vec n$ ein
  Normalenvektor und sei $Q$ ein beliebiger Punkt mit Ortsvektor $\vec q$.
  Dann gilt
  \[
    d(Q,E)
    =\left\lvert(\vec q-\vec p)\cdot\vec e_n\right\rvert
    =\frac{\left\lvert(\vec q-\vec p)\cdot\vec n\right\rvert}
           {\lvert\vec n\rvert},
    \qquad
    \vec e_n=\frac{\vec n}{\lvert\vec n\rvert}.
  \]
  Der Betrag macht den Abstand unabhängig von der Orientierung des
  Normalenvektors.
\end{TheoremBox}

\begin{ExampleBox}{Abstand eines Punktes von einer Ebene}{analytic_geometry_point_plane_distance_example}
  Gegeben sind
  \[
    E:\quad
    \vec x=\begin{pmatrix}1\\0\\0\end{pmatrix}
      +\lambda\begin{pmatrix}0\\1\\2\end{pmatrix}
      +\mu\begin{pmatrix}0\\2\\-1\end{pmatrix}
    \qquad\text{und}\qquad Q(3,2,2).
  \]
  Zunächst bestimmen wir einen Einheitsnormalenvektor:
  \begin{ClozeArea}[34]{11}
    \[
      \begin{pmatrix}0\\1\\2\end{pmatrix}
      \times
      \begin{pmatrix}0\\2\\-1\end{pmatrix}
      =\begin{pmatrix}-5\\0\\0\end{pmatrix}
      \quad\Longrightarrow\quad
      \vec e_n=\begin{pmatrix}1\\0\\0\end{pmatrix}.
    \]
  \end{ClozeArea}
  Damit folgt
  \begin{ClozeArea}[34]{7}
    \[
      d(Q,E)
      =\left\lvert
        \left(
          \begin{pmatrix}3\\2\\2\end{pmatrix}
          -\begin{pmatrix}1\\0\\0\end{pmatrix}
        \right)\cdot
        \begin{pmatrix}1\\0\\0\end{pmatrix}
      \right\rvert
      =2.
    \]
  \end{ClozeArea}
  Der Punkt $Q$ hat somit den Abstand $\Cloze{3cm}{2}$ von der Ebene.
\end{ExampleBox}
